% Minimal retained preamble for lane-p2 cards taken from
% sources/logic-fl05-ramsey-like-theorems-for-the-schre/pinned-original-source.tex
% (arXiv:2412.11598v3, CC BY 4.0).  This source ships no separate preamble file:
% the pinned tex is a concatenation of the authors' main.tex and macros.tex and
% carries its preamble inline, so the blocks below are sliced mechanically from
% that file and nothing has been renamed or redefined.
%
% Copied verbatim from the source:
%   line 2        : the authors' class line
%   lines 3-5,7,9 : graphicx, amsthm, the ams bundle, cleveref, url
%   line 10       : tikz-cd, used by the authors' commutative diagrams
%   lines 20-108  : the authors' macro definitions, taken in full so that every
%                   macro of the retained spans is available
%   lines 113-124 : the authors' theorem declarations
%   lines 131-141 : the authors' machinery for restating a theorem
%   lines 143-148 : the authors' coloured comment macros
%
% Omitted source preamble lines and why:
%   line 1, 110-111 and 150 (concatenation comment, title, date, author): the
%       assembler supplies the title block, author list and date;
%   line 6 (hyperref) and line 8 (inputenc with the utf8x option): the forced
%       xelatex engine supplies the hyperref wiring itself and utf8x is
%       incompatible with it;
%   line 115 (numberwithin theorem by section): the assembler numbers the
%       retained material itself;
%   lines 120 and 121-123 are declarations that the authors left commented out or
%       that duplicate earlier ones and are not reproduced.
\documentclass{article}
\usepackage{graphicx}
\usepackage{amsthm}
\usepackage{amsfonts, amsmath, amssymb, amsthm}
\usepackage{cleveref}
\usepackage{url}
\usepackage{tikz-cd}
\newcommand{\Psf}{\mathsf{P}}
\newcommand{\Qsf}{\mathsf{Q}}
\newcommand{\Rsf}{\mathsf{R}}
\def\F{\mathcal{F}}
\def\U{\mathcal{U}}
\def\C{\mathcal{C}}
\def\I{\mathcal{I}}
\def\D{\mathcal{D}}
\def\L{\mathcal{L}}
\def\V{\mathcal{V}}
\def\A{\mathcal{A}}
\def\B{\mathcal{B}}
\def\P{\mathcal{P}}
\def\M{\mathcal{M}}
\def\Nc{\mathcal{N}}
\def\Q{\mathcal{Q}}
\newcommand{\cs}{2^{\NN}}
\newcommand{\bs}{\NN^{\NN}}
\newcommand{\bstr}{2^{<\NN}}
\newcommand{\uh}{\upharpoonright}
\newcommand{\converge}{\!\!\downarrow}
\newcommand{\diverge}{\!\!\uparrow}
\newcommand{\qvdash}{\operatorname{{?}{\vdash}}}
\newcommand{\nqvdash}{\operatorname{{?}{\nvdash}}}
\newcommand{\ISig}{\mathsf{I}\Sigma^0}
\newcommand{\Robinson}{\mathsf{Q}}
\newcommand{\IPi}{\mathsf{I}\Pi^0}
\newcommand{\IDelta}{\mathsf{I}\Delta^0}
\newcommand{\Iopen}{\mathsf{I_{open}}}
\newcommand{\BSig}{\mathsf{B}\Sigma^0}
\newcommand{\BPi}{\mathsf{B}\Pi^0}
\newcommand{\BDelta}{\mathsf{B}\Delta^0}
\newcommand{\BCSig}{\mathsf{BC}\Sigma^0}
\newcommand{\BCPi}{\mathsf{BC}\Pi^0}
\newcommand{\BCDelta}{\mathsf{BC}\Delta^0}
\newcommand{\LSig}{\mathsf{L}\Sigma^0}
\newcommand{\LPi}{\mathsf{L}\Pi^0}
\newcommand{\CSig}{\mathsf{C}\Sigma}
\newcommand{\RCA}[0]{\mathsf{RCA}}
\newcommand{\WKL}[0]{\mathsf{WKL}}
\newcommand{\WWKL}[0]{\mathsf{WWKL}}
\newcommand{\ACA}[0]{\mathsf{ACA}}
\newcommand{\ATR}[0]{\mathsf{ATR}}
\newcommand{\PIOOCA}[0]{\Pi^1_1\mathsf{-CA}}
\newcommand{\CAC}[0]{\mathsf{CAC}}
\newcommand{\ADS}[0]{\mathsf{ADS}}
\newcommand{\RRT}[0]{\mathsf{RRT}}
\newcommand{\FS}[0]{\mathsf{FS}}
\newcommand{\PFS}[0]{\mathsf{PFS}}
\newcommand{\TS}[0]{\mathsf{TS}}
\newcommand{\GRRT}[0]{\mathsf{GRRT}}
\newcommand{\LRRT}[0]{\mathsf{LRRT}}
\newcommand{\EM}[0]{\mathsf{EM}}
\newcommand{\SEM}[0]{\mathsf{SEM}}
\newcommand{\HalfEM}[0]{\mathsf{Half}\mbox{-}\mathsf{EM}}
\newcommand{\RT}[0]{\mathsf{RT}}
\newcommand{\SRT}[0]{\mathsf{SRT}}
\newcommand{\NN}[0]{\mathbb{N}}
\newcommand{\nat}[0]{\mathbb{N}}
\newcommand{\card}{\operatorname{card}}
\newcommand{\finsub}{\subseteq_{\mathtt{fin}}}
\newcommand{\W}{\mathrm{W}}
\newcommand{\sW}{\mathrm{sW}}
\newcommand{\lex}{\mathrm{lex}}
\newcommand{\base}{\operatorname{base}}
\newcommand{\range}{\operatorname{ran}}
\def\qt#1{``#1''}%

\newtheorem*{statement}{Statement}
\newtheorem{theorem}{Theorem}
\numberwithin{theorem}{section}
\newtheorem{maintheorem}[theorem]{Main Theorem}
\newtheorem{lemma}[theorem]{Lemma}
\newtheorem{question}[theorem]{Question}
\newtheorem{proposition}[theorem]{Proposition}
%\newtheorem{statement}[theorem]{Statement}
\newtheorem{remark}[theorem]{Remark}
\newtheorem{definition}[theorem]{Definition}
\newtheorem{corollary}[theorem]{Corollary}
\newtheorem{example}[theorem]{Example}

% Allow to repeat theorems
\makeatletter
\newtheorem*{rep@theorem}{\rep@title}
\newcommand{\newreptheorem}[2]{%
\newenvironment{rep#1}[1]{%
 \def\rep@title{#2 \ref{##1}}%
 \begin{rep@theorem}}%
 {\end{rep@theorem}}}
\makeatother

\newreptheorem{theorem}{Theorem}

% Comments
\usepackage{xcolor}
\newcommand{\ludovic}[1]{\textcolor{blue}{#1}}
\newcommand{\quentin}[1]{\textcolor{green}{#1}}
\newcommand{\lorenzo}[1]{\textcolor{orange}{#1}}
\newcommand{\oriola}[1]{\textcolor{red}{#1}}
\title{A free finite block may be added to the stem while preserving property (b) of a condition}
\author{Lorenzo Carlucci, Oriola Gjetaj, Quentin Le Houérou, Ludovic Levy Patey}
\date{arxiv:2412.11598v3; CC BY 4.0}
\begin{document}
\maketitle
\noindent\textit{Editorial note.} The counted unit is the article's own Lemma lem:fs-preserving-b at source lines 1162-1167 together with its proof at source lines 1168-1175, a straight adjacent pairing in the pinned file. The pinned file is a concatenation of the authors' main and macro files; it is the article's fixed version arXiv:2412.11598v3 under CC BY 4.0.\par\medskip
\noindent\textit{Editorial note.} This source ships no separate preamble file, so the retained preamble is sliced mechanically from the pinned file by the lane and is shared, byte for byte, with the other cards of this source; it carries the authors' full macro block, including the restriction symbol that the retained statement uses.\par\medskip
\noindent\textit{Editorial note.} Three context spans carry the material the retained statement and proof rely on: the convention at line 1119 fixing the bounded-constrained colouring, the definition at line 1152 of the bounding function formed from the Catalan numbers, and the author's definition at lines 1154-1156 of the limit colouring attached to a stabilising reservoir. They are context, not counted units, and none overlaps the counted interval. The definition of a progressive function that the third case uses is recorded in the standard prerequisites instead of being carried as a span: that definition sits at line 1017, and any retained interval reaching back that far would also contain the still-open proof environment at line 1091 of a proposition that another lane counts, which the gate's interval self-check rightly refuses.\par\medskip
\noindent\textit{Editorial note.} The retained statement and proof contain no reference to another labelled result of the article and no citation, so this card carries no navigation pointer and no bibliography entry. In particular the two neighbouring lemmas of the same section are not named inside this span even though the section later combines all three.\par\medskip
\noindent\textit{Allowed standard prerequisites (editorial).}\begin{itemize}
\item Conditions, stabilisation, the limit colouring and progressiveness. The section works with conditions, each a triple of a colouring constrained by a computable bounding function, a finite stem and an infinite reservoir computing a fixed set D, with the two defining properties (a) and (b) of a condition; a reservoir stabilises an interval when the images of its finite sets land in fixed finite sets of short tuples, and to such a reservoir the section attaches a limit colouring whose values are the unions of those finite sets; a function into finite sets of natural numbers is progressive when every nonempty value has least element at least the least element of its argument. The author's own definitions are the allowed inputs.. The counted lemma is the core combinatorial step of the diagonalisation: it says exactly when a finite block of reservoir elements may be appended to the stem without destroying property (b), and the extensibility lemma later in the section combines it with two other lemmas of the same section.
\end{itemize}\medskip
In what follows, fix a coloring $f \in \P$ which is $b_f$-constrained, for some computable function $b_f : \NN \to \NN$.

Given a computable function $b : \NN \to \NN$, let $b ^{+} : \NN \to \NN$ be the computable function defined by  $b^{+}(m) = \sum_{n \leq h(m)} b(m) \times C_n$.

\begin{definition}
Let $X$ be a reservoir stabilizing $[0, k]$. The \emph{limit coloring} is the function $g_{k,X} : [k]^{\leq h(\cdot)} \to [k]^{<\omega}$ defined by $g_{k,X}(\emptyset) = \emptyset$ and $g_{k,X}(t) = \bigcup_{n \leq h(t) - |t|} I_{t, n}$ otherwise.
\end{definition}

\begin{lemma}\label[lemma]{lem:fs-preserving-b}
Let $(f, \sigma, X)$ be a condition
and $Y \subseteq X$ be an infinite set stabilizing $[0, k]$ for some~$k \in \NN$ and let $g_{k,Y} : [k]^{\leq h(\cdot)} \to [k]^{<\omega}$ be the limit coloring.
Let $\rho \subseteq X \uh_k$ be a finite $g_{k,Y}$-free set. Then $(f, \sigma \cup \rho, Y)$
satisfies property~(b).
\end{lemma}

\begin{proof}
Fix some~$s \in [\sigma \cup \rho \cup Y]^{\leq h(\cdot)}$. We have multiple cases.
\begin{itemize}
    \item Case 1: $s \cap \sigma \neq \emptyset$. Then by properties (a) and (b) of $(f, \sigma, X)$, $f(s) \cap (\sigma \cup X) \subseteq s$. Since $(\sigma \cup \rho) \subseteq (\sigma \cup X)$, then $f(s) \cap (\sigma \cup \rho) \subseteq s$.
    \item Case 2: $s \cap \sigma = \emptyset$ but $s \cap \rho \neq \emptyset$. By property (b) of $(f, \sigma, X)$, $f(s) \cap \sigma \subseteq s$. Let $t = s \cap \rho$ and $u = s \cap Y$. Since $\rho$ is $g_{k,Y}$-free, $g_{k,Y}(t) \cap \rho \subseteq t$. By definition of $g_{k,Y}(t)$, $I_{t, |u|} \subseteq g_{k,Y}(t)$, so $I_{t, |u|} \cap \rho \subseteq t$. In particular, $f(s) = f(t, u) \in I_{t, |u|}$, so $f(s) \cap \rho \subseteq t \subseteq s$. Thus, $f(s) \cap (\sigma \cup \rho) \subseteq s$.
    \item Case 3: $s \cap (\sigma \cup \rho) = \emptyset$. Then $s \subseteq Y$. Since $\min Y > \max (\sigma \cup \rho)$, then by progressiveness of~$f$, $f(s) \cap (\sigma \cup \rho) = \emptyset$.
\end{itemize}
\end{proof}

\end{document}
